\documentclass{ximera}

\input{../preamble.tex}


\title[Problem 8]{Problem 8}

\begin{document}
\begin{abstract} \end{abstract}
\maketitle


% Extracted from approximatingTheAreaUnderACurve.tex, problem #8
\begin{problem}
We want to approximate the area under the curve  using a right Riemann Sum with the given value of $n$. Write the sum in summation notation and evaluate it.
	\begin{enumerate}
	%part a
	\item  $\sin (x)$, \; $\left[ 0, \dfrac{\pi}{2} \right]$, \; $n=3$
	
		\begin{explanation}
		
		$\Delta x=\dfrac{b-a}{n}=\dfrac{\dfrac{\pi }{2}-0}{3}=\dfrac{\pi }{6}$.  \\
		
		$x_k^*=x_k=a+k\Delta x=0+\dfrac{\pi }{6}k=\dfrac{\pi }{6}k.$
		\begin{align*}
		\displaystyle \sum_{k=1}^{3}  f ( x_k^* ) \Delta x  &= \displaystyle \sum_{k=1}^{3} \left( f \left( \dfrac{\pi}{6}k \right) \cdot \dfrac{\pi}{6} \right) \\
		&= \dfrac{\pi}{6} \displaystyle \sum_{k=1}^{3} \sin \left( \dfrac{\pi}{6} k \right) \\
		&= \dfrac{\pi}{6} \left( \sin \left( \dfrac{\pi}{6} \right) + \sin \left( \dfrac{2 \pi}{6} \right) + \sin \left( \dfrac{3 \pi}{6} \right) \right) \\
		&= \dfrac{\pi}{6} \left( \sin \left( \dfrac{\pi}{6} \right) + \sin \left( \dfrac{\pi}{3} \right) + \sin \left( \dfrac{\pi}{2} \right) \right) \\
		&= \dfrac{\pi}{6} \left( \dfrac{1}{2} + \dfrac{\sqrt{3}}{2} + 1 \right)  \\
		&= \dfrac{\pi}{6} \left( \dfrac{3}{2} + \dfrac{\sqrt{3}}{2} \right)  \\
		&= \dfrac{\pi}{12} (3 + \sqrt{3} )
		\end{align*}
		\end{explanation}
\item  $f(x) = x^2 - 9x + 18$, \; $[7,10]$, \; $n=6$
		\begin{explanation}
		$\Delta x=\dfrac{b-a}{n}=\dfrac{10-7}{6}=\dfrac{1}{2}.$  \\
		$x_k^* =x_k=a+k\Delta x=7+\dfrac{1}{2}k.$  
		\begin{align*}
		\displaystyle \sum_{k=1}^{6} f(x_k^*) \Delta x &=f(7+\dfrac{1}{2})\cdot\dfrac{1}{2}+f(7+1)\cdot\dfrac{1}{2}+f(7+\dfrac{3}{2})\cdot\dfrac{1}{2}+f(7+2)\cdot\dfrac{1}{2}+f(7+\dfrac{5}{2})\cdot\dfrac{1}{2}+f(7+3)\cdot\dfrac{1}{2}\\
		&=f(7.5)\cdot\dfrac{1}{2}+f(8)\cdot\dfrac{1}{2}+f(8.5)\cdot\dfrac{1}{2}+f(9)\cdot\dfrac{1}{2}+f(9.5)\cdot\dfrac{1}{2}+f(10)\cdot\dfrac{1}{2}\\
&=\dfrac{1}{2}\left(f(7.5)+f(8)+f(8.5)+f(9)+f(9.5)+f(10)\right)\\
&=\dfrac{1}{2}\left((7.5^2-9(7.5)+18)+(8^2-9(8)+18)+(8.5^2-9(8.5)+18)+(9^2-9(9)+18)+(9.5^2-9(9.5)+18)+(10^2-9(10)+18)\right)\\
&=\dfrac{1}{2}\left((7.5^2+8^2+8.5^2+9^2+9.5^2+10^2)-9(7.5+8+8.5+9+9.5+10)+6(18)\right)\\
&= \dfrac{397}{8}
		\end{align*}
		Or in sigma notation:
		\begin{align*}
		\displaystyle \sum_{k=1}^{6} f(a + k \Delta x) \Delta x &= \displaystyle \sum_{k=1}^{6} \left( f \left( 7 + \dfrac{1}{2} k \right) \cdot \dfrac{1}{2} \right)  \\
		&= \dfrac{1}{2} \displaystyle \sum_{k=1}^{6} \left( \left( 7 + \dfrac{1}{2} k \right)^2 - 9 \left( 7 + \dfrac{1}{2} k \right) + 18 \right) \\
		&= \dfrac{1}{2} \displaystyle \sum_{k=1}^{6} \left( 49 + 7k + \dfrac{1}{4} k^2 - 63 - \dfrac{9}{2} k + 18 \right) \\
		&= \dfrac{1}{2} \displaystyle \sum_{k=1}^{6} \left( 4 + \dfrac{5}{2} k+ \dfrac{1}{4} k^2 \right) \\
		&= \dfrac{1}{2} \left( 4 \displaystyle \sum_{k=1}^{6} 1 + \dfrac{5}{2} \displaystyle \sum_{k=1}^{6} k+ \dfrac{1}{4} \displaystyle \sum_{k=1}^{6} k^2 \right) \\
		&= \dfrac{1}{2} \left( 4(6) + \dfrac{5}{2} \cdot \dfrac{(6)(7)}{2} + \dfrac{1}{4} \cdot \dfrac{(6)(7)(13)}{6} \right) \\
		&= \dfrac{1}{2} \left( 24 + \dfrac{105}{2} + \dfrac{91}{4} \right) \\
		&= \dfrac{1}{2} \cdot \dfrac{96 + 210 + 91}{4} = \dfrac{397}{8}
		\end{align*}
		\end{explanation}
%part b
	
		
		
		
	\end{enumerate}
		
		
		

\end{problem}



\end{document}
